\section*{Description of test files to be used for evaluating the respective milestone}

The main evaluation artifacts for this milestone are the automated test modules, the shared test infrastructure, and the \texttt{.tl} fixture programs.

\subsection*{Automated test modules}

\begin{longtable}{p{0.22\textwidth} p{0.72\textwidth}}
\toprule
\textbf{File} & \textbf{Purpose} \\
\midrule
\texttt{tests/test\_default.py} & Checks properties in \texttt{default} mode, where the entire state starts from a concrete zero initialization. These tests focus on whether the verifier can prove or refute invariants from a fixed start state. \\
\texttt{tests/test\_universal.py} & Checks properties in \texttt{universal} mode, where user variables and key turtle components are unconstrained initially. These tests validate reasoning under arbitrary initial conditions. \\
\texttt{tests/test\_specific.py} & Checks properties in \texttt{specific} mode using explicit parameter dictionaries. These tests validate the part of the pipeline that reads concrete input assignments from the command line and injects them into the start state. \\
\texttt{tests/test\_safety\_range.py} & Checks properties in \texttt{safety-range} mode. These tests validate the ghost-parameter mechanism and check whether the implementation can recover a symbolic safe region over initial inputs. \\
\texttt{tests/helpers.py} & Shared testing infrastructure. It builds the fixed-point system for a given program, constructs the property-evaluation context, suppresses verbose pipeline output during testing, and provides helper assertions for expected pass/fail outcomes. \\
\texttt{tests/run\_all.py} & Top-level test runner that discovers all \texttt{test\_*.py} files and executes the complete suite. \\
\bottomrule
\end{longtable}

\subsection*{Fixture programs used by the tests}

\begin{longtable}{p{0.28\textwidth} p{0.66\textwidth}}
\toprule
\textbf{Program file} & \textbf{Purpose in evaluation} \\
\midrule
\texttt{assign\_basic.tl} & Basic sequential assignments; used to test simple arithmetic invariants and no-motion/no-pen behavior. \\
\texttt{assign\_algebra.tl} & Multi-step arithmetic including multiplication and division; used to test algebraic propagation and nonlinear cases. \\
\texttt{conditional.tl} & Simple conditional branch; used to test branching semantics. \\
\texttt{loop\_basic.tl} & Counted loop with accumulation; used to test loop invariants. \\
\texttt{loop\_accum.tl} & Repetition-driven accumulator growth; used to test monotonic arithmetic invariants. \\
\texttt{loop\_cond.tl} & Loop with internal branching and arithmetic updates; used to test interaction between repetition and conditionals. \\
\texttt{goto\_computed.tl} & Computed \texttt{goto} destinations from symbolic variables; used for geometric safety properties. \\
\texttt{square\_goto.tl} & Multiple \texttt{goto} commands that trace a square; used to test bounded-region properties. \\
\texttt{forward\_square.tl} & Motion using \texttt{forward} and \texttt{right}; used for trig-dependent geometric tests. \\
\texttt{turns\_only.tl} & Pure heading updates; used for directional reasoning and normalization behavior. \\
\texttt{pen\_only.tl} & Pen-state updates without arithmetic or motion; used to validate \texttt{pendown}/\texttt{penup} semantics. \\
\texttt{pen\_with\_var.tl} & Pen commands combined with variable updates; used to test mixed symbolic and boolean state. \\
\texttt{pen\_toggle.tl} & Alternating pen state changes; useful for repeated pen-state transitions. \\
\texttt{param\_goto.tl} & Parameterized \texttt{goto}; used in \texttt{specific} and \texttt{safety-range} modes. \\
\texttt{param\_loop.tl} & Loop whose initial value depends on a parameter; used to test parameter-sensitive arithmetic invariants. \\
\texttt{param\_cond.tl} & Branching behavior controlled by a parameter; used to test mode-dependent initialization and branch selection. \\
\texttt{param\_scale.tl} & Arithmetic scaling with an input parameter; used to test concrete parameter injection and sign-sensitive properties. \\
\texttt{param\_pen.tl} & Parameterized arithmetic followed by pen commands; used to test the interaction of user inputs and pen state. \\
\texttt{read\_before\_write.tl} & Reads an input parameter before any local assignment; used to validate handling of externally supplied variables in \texttt{specific} mode. \\
\bottomrule
\end{longtable}

\subsection*{How the test files should be used during evaluation}

For milestone evaluation, the recommended procedure is:
\begin{enumerate}[leftmargin=1.5em]
    \item run the automated suite with \texttt{python run\_all.py};
    \item inspect individual mode-specific modules when needed; and
    \item use representative fixture programs directly with \texttt{safety\_properties.py} to demonstrate interactive verification in \texttt{default}, \texttt{universal}, \texttt{specific}, and \texttt{safety-range} modes.
\end{enumerate}

This combination of direct verifier runs and automated regression tests provides a practical and reproducible way to evaluate the milestone.


\subsubsection*{Automated validation through unit tests}
The repository contains a dedicated automated test suite under \texttt{Chiron-Framework/ChironCore/CHC\_Verification/tests}. From code inspection, the suite currently contains 89 unit tests spanning the four supported modes. The tests are organized around arithmetic, geometric, pen-state, directional, and safety-range behavior, and they collectively check both expected invariants and expected violations. A shared helper layer suppresses pipeline debug output, constructs the fixed-point system, creates the evaluation context for properties, and provides convenience assertions for pass/fail verdicts.

The full automated suite can be run as follows:

\begin{verbatim}
cd Chiron-Framework/ChironCore/CHC_Verification/tests
python run_all.py
\end{verbatim}

Useful narrower invocations are also supported:

\begin{verbatim}
# verbose test run
python run_all.py -v

# run one test module
python -m unittest test_default

# run one class
python -m unittest test_specific.TestSpecificArithmetic

# run one individual test
python -m unittest \
  test_default.TestDefaultArithmetic.test_arith_assign_z_nonneg_pass
\end{verbatim}

Overall, the work completed so far establishes a usable verification tool rather than just an isolated encoding: there is a state extractor, a mode-sensitive invariant generator, a step-rule translator, a property interface, and an automated regression suite.
